% =====================================================================
%  Buhera OS Module Integration:
%  A Contract, Reference, and Construction Manual for New Modules
%  in the Buhera Federation
%
%  Self-contained blueprint. Any collaborator (human or AI) reading this
%  document alone should be able to produce a working Buhera module that
%  slots into the federation, without reading the surrounding codebase.
%  Every path, every trait, every wiring step, and every failure mode is
%  named here.
% =====================================================================
\documentclass[11pt,a4paper]{article}

\usepackage[utf8]{inputenc}
\usepackage[T1]{fontenc}
\usepackage{lmodern}
\usepackage[a4paper,margin=2.3cm]{geometry}
\usepackage{amsmath,amssymb,amsthm,mathtools}
\usepackage{bm}
\usepackage{booktabs}
\usepackage{array}
\usepackage{enumitem}
\usepackage{microtype}
\usepackage{listings}
\usepackage{xcolor}
\usepackage{tabularx}
\usepackage[colorlinks=true,linkcolor=blue!45!black,
            citecolor=blue!45!black,urlcolor=blue!45!black]{hyperref}
\usepackage{cleveref}

% ---- Listings style ----
\definecolor{lstkw}{RGB}{30,64,175}
\definecolor{lstcomment}{RGB}{22,101,52}
\definecolor{lststring}{RGB}{159,18,57}
\definecolor{lstnum}{RGB}{180,83,9}
\definecolor{lstbg}{RGB}{250,250,250}

\lstdefinestyle{jscode}{
  language=JavaScript,
  basicstyle=\ttfamily\footnotesize,
  keywordstyle=\color{lstkw}\bfseries,
  commentstyle=\color{lstcomment}\itshape,
  stringstyle=\color{lststring},
  numberstyle=\color{lstnum},
  backgroundcolor=\color{lstbg},
  frame=single,framesep=4pt,xleftmargin=6pt,xrightmargin=4pt,
  breaklines=true,showstringspaces=false,columns=fullflexible,tabsize=2,
  morekeywords={async,await,export,import,from,const,let,return,typeof,
    Promise,Array,Set,Map,Number,String,Boolean,Object,JSON,fetch}
}
\lstdefinestyle{shell}{
  basicstyle=\ttfamily\footnotesize,
  backgroundcolor=\color{lstbg},
  frame=single,framesep=4pt,xleftmargin=6pt,xrightmargin=4pt,
  breaklines=true,columns=fullflexible,tabsize=2
}
\lstset{style=jscode}

% ---- Theorem-like environments repurposed as design blocks ----
\newtheoremstyle{contract}{\topsep}{\topsep}{}{}{\bfseries}{.}{.5em}{}
\theoremstyle{contract}
\newtheorem{contract}{Contract}[section]
\newtheorem{invariant}[contract]{Invariant}
\newtheorem{recipe}[contract]{Recipe}
\newtheorem{caution}[contract]{Caution}

\title{\textbf{Buhera OS Module Integration:\\
A Contract, Reference, and Construction Manual\\
for New Modules in the Buhera Federation}}

\author{%
Kundai Farai Sachikonye\\
\small Technical University of Munich\\
\small AIMe Registry for Artificial Intelligence\\
\small \texttt{kundai.sachikonye@wzw.tum.de}}

\date{\today}

\begin{document}
\maketitle

\begin{abstract}
\noindent
Buhera OS is a small operating system for finite reasoning agents. It
runs entirely in the browser as a Next.js application (single-page,
client-side after mount), presents a terminal, and dispatches
instructions to a federation of registered modules. Each module owns
one domain (memory, mass spectrometry, natural-language interception,
individuation-theoretic search, tandem context carry, LLM synthesis,
and so on). The federation grows by addition, never by modification of
existing modules. This document specifies exactly what a new module
must implement, what it may assume, and how it is wired into the running
system. It is a contract document, not a survey: every section states
what MUST hold, and where a design choice is genuinely open it says so
explicitly. A collaborator reading this document alone (with no prior
exposure to the Buhera codebase) should be able to produce a compiling,
registerable, dispatchable module that appears in \texttt{:modules},
returns valid \texttt{ActResult}s, and interoperates with the existing
federation without altering a single line of the shipped modules.
\end{abstract}

\noindent\textbf{Scope.} New scientific/domain modules for the Buhera OS
frontend at \texttt{long-grass/}. The scheduler, the substrate, the
kernel, the audit log, and the LLM cascade are already implemented and
are not part of this document except at the interface where new modules
touch them.

\tableofcontents

\vspace{6pt}\hrule\vspace{6pt}

% =====================================================================
\part{The System a Module Is Joining}
% =====================================================================

% =====================================================================
\section{Architectural Overview}
\label{sec:overview}
% =====================================================================

Before writing a line of module code, a collaborator must know four
things about the system the module joins.

\subsection{The runtime}

Buhera OS is a Next.js 13 application (Pages Router) served from
\texttt{long-grass/}. The terminal is a single React component,
\texttt{src/components/BuheraTerminal.js}, mounted at the root route.
After hydration the entire federation runs client-side; the browser is
the operating system. Server-side API routes exist only to hold
credentials (LLM keys) that must not appear in the client bundle.

The deployment target is Vercel. The build is \texttt{next build}; the
dev loop is \texttt{next dev}. There is no separate installation step
for a user: opening the deployed URL is running the OS.

\subsection{The three roles a module can occupy}

Every existing Buhera module falls into one of three functional roles,
and every new module MUST declare which role(s) it occupies:

\begin{enumerate}[leftmargin=1.6em,itemsep=2pt]
\item \textbf{Solver.} Given an instruction, do the work and return an
answer. Examples: \texttt{lavoisier} (forward-simulate mass spectra),
\texttt{vahera} (execute a memory DSL script), \texttt{graffiti}
(execute a \texttt{.grf} search program).
\item \textbf{Router / Interceptor.} Given natural-language input,
translate to a routable instruction. Example: \texttt{zangalewa} (MSI --
utterance to S-coord + rendered card).
\item \textbf{Meta / Context.} Read the accumulated audit log or session
state and return derived information. Example: \texttt{purpose-carry}
(tandem carry over accumulated steps).
\end{enumerate}

A module MAY occupy more than one role. A module MAY NOT invent a
fourth role without first extending this document.

\subsection{The three surfaces a module touches}

Independently of role, every module touches at most three surfaces:

\begin{enumerate}[leftmargin=1.6em,itemsep=2pt]
\item \textbf{The registry} (\S\ref{sec:registry}) --- always. Every
module is registered and dispatched through the same trait.
\item \textbf{The terminal renderer} (\S\ref{sec:renderer}) --- if the
module returns output that should display specially (a chart, a card, a
table). Purely textual output uses the built-in \texttt{ArtifactText}
and requires no renderer.
\item \textbf{Server-side API routes} (\S\ref{sec:server}) --- if the
module needs credentials, network access, or Node-only APIs
(filesystem, LLM SDKs). A module that runs entirely in the browser
skips this surface.
\end{enumerate}

Anything outside these three surfaces is the module's private business.
The registry does not care what the module does internally.

\subsection{Naming and identity}

Every module has a stable string \texttt{id}. Ids are lowercase,
kebab-cased if compound, never namespaced. Good: \texttt{lavoisier},
\texttt{purpose-carry}, \texttt{crispr-designer}. Bad:
\texttt{@buhera/crispr}, \texttt{CRISPRDesigner}, \texttt{crispr\_v2}.

The \texttt{id} appears in the audit log, in \texttt{:modules}, in
turbulance scripts as \texttt{dispatch("<id>", ...)}, and in the
module's error messages. Once shipped, the \texttt{id} is a permanent
identifier; renaming it is a breaking change to every caller.

% =====================================================================
\section{The Module Trait}
\label{sec:registry}
% =====================================================================

The trait is defined in \texttt{src/lib/modules/registry.js}. Every
module conforms to it verbatim; the registry is not extensible on this
axis. This section is the load-bearing specification of the whole
document: if a module implements this trait correctly, everything else
in the federation just works.

\subsection{The trait}

\begin{lstlisting}
export const myModule = {
  id: "my-module",

  describe() { /* returns metadata; see 2.2 */ },

  async execute(instruction, actBudget = 1) {
    /* the module's work; returns an ActResult; see 2.3 */
  },

  outputCell(instruction) {
    /* optional sufficiency descriptor; see 2.4 */
    return { kind: "my_cell" };
  },
};
\end{lstlisting}

Four fields. The \texttt{id} is a string. The \texttt{describe},
\texttt{execute}, and \texttt{outputCell} are functions. Only
\texttt{execute} may be \texttt{async}; \texttt{describe} and
\texttt{outputCell} MUST be synchronous.

\subsection{describe()}

Returns a plain object with three fields:

\begin{lstlisting}
describe() {
  return {
    id: "my-module",
    description: "one-sentence purpose of the module",
    instructions: [
      'dispatch("my-module", "example simple invocation")',
      'dispatch("my-module", { kind: "action", param: value })',
    ],
  };
}
\end{lstlisting}

\begin{contract}[describe]
\label{contract:describe}
\texttt{describe()} MUST be pure (no side effects), synchronous, and
inexpensive. The terminal calls it on every render of \texttt{:modules}.
It MUST return an object with the exact three fields shown. The
\texttt{instructions} array carries example invocations shown to the
user; it should demonstrate the module's shape, not enumerate every
possible instruction.
\end{contract}

\subsection{execute(instruction, actBudget)}

The core method. Called by the registry when a caller dispatches to
this module.

\subsubsection*{The \texttt{instruction} argument}

Whatever the caller passed as the second argument of
\texttt{dispatch()}. It MAY be a string, an object, a number, or
\texttt{null}. A module MUST accept a bare string as its most common
form (interactive terminal callers type strings). Modules that require
structured input MUST accept a plain string as a request for a default
demo, or produce a clear error explaining what shape is expected.

The convention across existing modules: an instruction of \texttt{""}
or \texttt{"demo"} runs a canned demonstration. This is not required
but is strongly recommended for discoverability.

\subsubsection*{The \texttt{actBudget} argument}

A hint about thoroughness. Defaults to \texttt{1}. Larger values indicate
the caller wants the module to invest more effort. Modules that are
deterministic and non-iterative may ignore this argument; modules that
have a knob for depth/quality/iterations SHOULD honour it.

\subsubsection*{The return value: \texttt{ActResult}}

Every call to \texttt{execute} MUST return a Promise resolving to an
object of the following shape:

\begin{lstlisting}
{
  ok: boolean,
  output_delta: {
    kind: string,
    // ...any other fields the renderer for this kind expects
  },
  residue: number,   // >= 0
  completed: boolean,
  error?: string,    // present iff ok === false
}
\end{lstlisting}

\begin{contract}[ActResult]
\label{contract:actresult}
Every field is mandatory unless marked optional. The registry does not
apply defaults; a module returning \texttt{\{value: 42\}} instead of an
ActResult is broken.

\begin{itemize}[leftmargin=1.4em,itemsep=1pt]
\item \texttt{ok} is a boolean success indicator. \texttt{false} means
the module received a valid dispatch but could not complete it (missing
input, upstream failure, invalid state). It does not mean the module
crashed --- crashes throw from the caller's perspective, and the
registry converts them to \texttt{ok:false} automatically.
\item \texttt{output\_delta.kind} is a string tag naming the shape of
the rest of the output. It selects which renderer runs. If no custom
renderer is registered (\S\ref{sec:renderer}), \texttt{kind: "text"}
with \texttt{lines: [string]} produces a plain-text fallback.
\item \texttt{residue} is a non-negative number the caller may use as a
progress or content signal. See \S\ref{sec:residue} for what to put
here.
\item \texttt{completed} is \texttt{true} when the module has nothing
more to produce for this instruction. All current modules return
\texttt{true} unconditionally; streaming or partial-result support is a
future extension.
\item \texttt{error} is a human-readable error string. It MUST be
present when \texttt{ok} is \texttt{false} and absent when \texttt{ok}
is \texttt{true}.
\end{itemize}
\end{contract}

\begin{caution}[Do not throw]
\label{caution:throw}
A module SHOULD NOT let exceptions propagate out of \texttt{execute}.
The registry does catch them and converts them to a generic
\texttt{ok:false} ActResult, but the error message will be less specific
than one the module produces itself. Wrap the module's work in a
\texttt{try/catch} at the top of \texttt{execute} and return a proper
ActResult with a domain-specific error string.
\end{caution}

\subsection{outputCell(instruction)}

Returns a descriptor of what kind of output the module would produce for
this instruction, without producing it. This is used by future
sufficiency-check machinery (the scheduler's gate) to decide whether a
dispatch is worth performing. In the current federation
\texttt{outputCell} is a stub returning \texttt{\{kind: "<something>\_cell"\}};
a real implementation will land when the scheduler needs it.

\begin{contract}[outputCell]
\label{contract:outputcell}
Every module MUST implement \texttt{outputCell}. It MUST be synchronous
and cheap. A stub returning \texttt{\{kind: "<module\_id>\_cell"\}} is
acceptable and is what current modules use.
\end{contract}

% =====================================================================
\section{The Registry, the Audit Log, and Post-Dispatch Hooks}
\label{sec:registry-mechanics}
% =====================================================================

Every dispatched act flows through the same code path in
\texttt{src/lib/modules/registry.js}. A module writer does not modify
this file; understanding it is nonetheless necessary because the
mechanics affect what a module can rely on.

\subsection{The dispatch cycle}

The registry's \texttt{dispatch(moduleId, instruction, actBudget)}
performs, in order:

\begin{enumerate}[leftmargin=1.6em,itemsep=1pt]
\item Look up the module by \texttt{moduleId}; throw if unknown.
\item Record the wall-clock start time.
\item Call \texttt{module.execute(instruction, actBudget)} in a
\texttt{try/catch}. Exceptions are converted to a synthetic
\texttt{ok:false} ActResult.
\item Construct an audit-log entry: monotone \texttt{act\_id}, module
id, verbatim instruction, act budget, the returned result, wall-clock
duration, ISO timestamp.
\item Push the entry onto the append-only audit log.
\item Run every registered post-dispatch hook against the entry.
Hooks are isolated: an exception in one hook is logged and swallowed.
\item Return the ActResult to the caller.
\end{enumerate}

The audit log is in-memory. It resets on page reload. A future
persistence layer will not require modules to change.

\subsection{Post-dispatch hooks}

Hooks let observers watch every dispatched act without the modules
knowing about them. The signature is \texttt{(entry) => void}. Hooks
are registered by consumers (currently: the purpose-carry session
feeder in \texttt{BuheraTerminal.js}); modules do not register their
own hooks.

\begin{invariant}[Modules are hook-oblivious]
\label{inv:hookoblivious}
A module MUST NOT rely on any specific hook being registered, nor on
its execution order. Modules interact with the world only through their
\texttt{execute()} return value.
\end{invariant}

\subsection{What a module may read from the registry}

A module may import and call:

\begin{itemize}[nosep,leftmargin=1.4em]
\item \texttt{getAuditLog()} --- returns a shallow copy of the audit
log; useful for meta modules like purpose-carry that reason about
history.
\item \texttt{listModules()} --- returns the metadata of every
registered module; useful for a routing/interceptor module.
\item \texttt{getModule(id)} --- returns another module's trait object;
rarely needed.
\item \texttt{dispatch(id, instr, actBudget)} --- lets one module
dispatch to another. Use sparingly; deep call chains are hard to debug.
\end{itemize}

A module MUST NOT import \texttt{clearAuditLog}, \texttt{register}, or
\texttt{onDispatch}. Those belong to the terminal's bootstrap.

% =====================================================================
\section{Terminal Wiring}
\label{sec:renderer}
% =====================================================================

Two edits to \texttt{src/components/BuheraTerminal.js} land every new
module: registration and (optionally) a renderer. Nothing else in the
file changes.

\subsection{Registration}

At the top of the file, add a new import next to the existing module
imports:

\begin{lstlisting}
import { myModule } from "@/lib/modules/my-module";
\end{lstlisting}

Inside the mount \texttt{useEffect}, add one \texttt{register()} call:

\begin{lstlisting}
useEffect(() => {
  kernelRef.current = bootBlank();
  register(vaheraModule);
  register(echoModule);
  register(lavoisierModule);
  register(purposeModule);
  register(zangalewaModule);
  register(graffitiModule);
  register(purposeCarryModule);
  register(myModule);           // <-- new
  // ...post-dispatch hooks below
}, []);
\end{lstlisting}

That is the only file change required for a module that returns plain
text.

\subsection{Renderers}

A module producing rich output (a chart, a card, a table) SHOULD supply
a React component and register a case in the artifact switch. The
convention is:

\begin{lstlisting}
function ArtifactMyModule({ ...fields }) {
  return (
    <div className="text-gray-300">
      {/* render the fields */}
    </div>
  );
}
\end{lstlisting}

Then in the \texttt{Artifact} switch, add a case:

\begin{lstlisting}
case "my_module_result":
  return <ArtifactMyModule
    fieldA={result.fieldA}
    fieldB={result.fieldB}
  />;
\end{lstlisting}

The \texttt{result} argument here is the \texttt{output\_delta} object
the module returned. The switch is keyed on \texttt{result.kind}, so
the \texttt{output\_delta.kind} string chosen by the module MUST match
the \texttt{case} label.

\begin{recipe}[Naming a kind]
\label{recipe:kind}
Use \texttt{<module\_id>\_<result\_type>}. Examples from the current
federation: \texttt{lavoisier\_run}, \texttt{graffiti\_result},
\texttt{zangalewa\_render}, \texttt{purpose\_synthesis},
\texttt{purpose\_carry}, \texttt{purpose\_carry\_stats}. Consistent
naming makes the switch statement legible and prevents accidental
collisions.
\end{recipe}

\subsection{Textual output as fallback}

A module that returns

\begin{lstlisting}
{
  ok: true,
  output_delta: {
    kind: "text",
    lines: ["first line", "second line", ...]
  },
  residue: 1,
  completed: true,
}
\end{lstlisting}

renders as multi-line grey text with no custom code required. This is
the fastest path to a working module and a legitimate first step for a
new module before a proper renderer is designed.

% =====================================================================
\section{Server-Side API Routes}
\label{sec:server}
% =====================================================================

A module that needs credentials, filesystem access, or Node-only APIs
splits into two files:

\begin{enumerate}[nosep,leftmargin=1.4em]
\item A client-side adapter under \texttt{src/lib/modules/} that
conforms to the Module trait. Its \texttt{execute} makes a
\texttt{fetch()} call to a server route.
\item A server route under \texttt{src/pages/api/} that does the
actual work and returns JSON.
\end{enumerate}

This split is not optional for modules touching credentials --- the
browser bundle MUST NOT contain any secret. Vercel's environment
variables reach only the server side.

\subsection{The server-route contract}

An API route is a plain Next.js API handler:

\begin{lstlisting}
export default async function handler(req, res) {
  if (req.method !== "POST") {
    res.setHeader("Allow", "POST");
    return res.status(405).json({ ok: false, error: "method not allowed" });
  }

  const { field1, field2 } = req.body || {};
  if (typeof field1 !== "string") {
    return res.status(400).json({ ok: false, error: "field1 required" });
  }

  try {
    const result = await doTheWork(field1, field2);
    return res.status(200).json({ ok: true, ...result });
  } catch (err) {
    return res.status(502).json({
      ok: false,
      error: err.message || String(err),
    });
  }
}
\end{lstlisting}

Conventions:

\begin{itemize}[nosep,leftmargin=1.4em]
\item POST only. GET is reserved for introspection endpoints like
\texttt{/api/providers}.
\item JSON body in, JSON body out.
\item Errors return a JSON body with \texttt{ok: false} and a
descriptive \texttt{error} string; the HTTP status distinguishes the
failure class (400 bad input, 500 our fault, 502 upstream fault, 503
missing configuration).
\item Never throw across the handler boundary. Every error path returns
a JSON body.
\end{itemize}

\subsection{The client-side adapter}

The adapter posts to the route and translates the response into an
ActResult:

\begin{lstlisting}
export const myModule = {
  id: "my-module",
  describe() { /* ... */ },

  async execute(instruction, _actBudget = 1) {
    const body = normaliseInstruction(instruction);

    let res;
    try {
      res = await fetch("/api/my-module", {
        method: "POST",
        headers: { "Content-Type": "application/json" },
        body: JSON.stringify(body),
      });
    } catch (err) {
      return {
        ok: false,
        output_delta: {
          kind: "text",
          lines: [`my-module: network error — ${err.message}`],
        },
        residue: 0,
        completed: true,
        error: err.message,
      };
    }

    const json = await res.json().catch(() => null);
    if (!json || !json.ok) {
      return {
        ok: false,
        output_delta: {
          kind: "text",
          lines: [`my-module: ${json?.error ?? `HTTP ${res.status}`}`],
        },
        residue: 0,
        completed: true,
        error: json?.error ?? `HTTP ${res.status}`,
      };
    }

    return {
      ok: true,
      output_delta: {
        kind: "my_module_result",
        ...json,
      },
      residue: 1,
      completed: true,
    };
  },

  outputCell() { return { kind: "my_module_cell" }; },
};
\end{lstlisting}

\subsection{Using the shared LLM cascade}

Modules that need LLM output SHOULD use the shared cascade in
\texttt{src/lib/server/llm-cascade.js} rather than calling providers
directly. The cascade tries Ollama, then Gemini, then HuggingFace, then
OpenAI, using whichever is configured. Signature:

\begin{lstlisting}
import { chatCascade } from "@/lib/server/llm-cascade";

const result = await chatCascade({
  system: "system prompt...",
  user: "user prompt...",
  jsonSchema: { /* optional; requests structured output */ },
  maxTokens: 2048,
  temperature: 0.2,
  only: ["gemini"], // optional: restrict to a subset
});

if (!result.ok) {
  // result.error, result.stage, result.provider
}
// result.content: string
// result.provider: which provider actually answered
// result.model: which model within that provider
\end{lstlisting}

\begin{contract}[Cascade usage]
\label{contract:cascade}
Modules SHOULD NOT call OpenAI, Gemini, HuggingFace, or Ollama APIs
directly. All LLM calls flow through \texttt{chatCascade}. This
centralises the ``which key is set'' logic, the fallback logic, and the
JSON-schema translation for structured output. A module that ships a
direct provider call will silently break when the deployment switches
providers.
\end{contract}

% =====================================================================
\section{The Residue Field: What to Put There}
\label{sec:residue}
% =====================================================================

Every ActResult carries a \texttt{residue} number. Its semantics are set
by the scheduling-mechanism paper and the tandem-carry paper: residue is
the measured distance of a partial result from a recognised solution,
with a strictly positive floor $\beta > 0$.

For a new module, pick a residue interpretation that is:

\begin{itemize}[nosep,leftmargin=1.4em]
\item \textbf{Non-negative} --- always.
\item \textbf{Monotone in content produced} --- more work done should
correspond to a specific direction (higher or lower), consistently
across every dispatch of the module.
\item \textbf{Interpretable by callers} --- the residue is a signal, so
picking arbitrary numbers is worse than picking any consistent rule.
\end{itemize}

Common patterns:

\begin{itemize}[nosep,leftmargin=1.4em]
\item \textbf{Content-count residue.} \texttt{residue = records.length}
(lavoisier), \texttt{= number of yielded claims} (graffiti). Higher =
more produced.
\item \textbf{Confidence proxy.} \texttt{residue = 1 - confidence}
(zangalewa's \texttt{S\_e} approximation). Lower = closer to done.
\item \textbf{Trivial signal.} \texttt{residue = 0} for a pure query,
\texttt{residue = 1} for a successful action. Acceptable for modules
whose progress isn't naturally continuous.
\end{itemize}

The residue field is not currently consumed by any scheduler in the
frontend; the scheduler in \texttt{src/lib/scheduler/} reads it in
principle but is not yet wired into the dispatch path. Nevertheless,
modules ship with valid residues so that when the wiring lands, the
module works without change.

% =====================================================================
\part{Building a New Module}
% =====================================================================

% =====================================================================
\section{Directory Layout and File Skeleton}
\label{sec:skeleton}
% =====================================================================

\subsection{Where files live}

\begin{tabularx}{\linewidth}{lX}
\toprule
\textbf{Path} & \textbf{What lives there} \\
\midrule
\texttt{src/lib/modules/<id>-module.js} & The Module trait implementation. Required. \\
\texttt{src/lib/<domain>/} & Domain library code (helpers, algorithms, data). Optional. \\
\texttt{src/pages/api/<route>.js} & Server routes. Optional. \\
\texttt{src/lib/server/} & Server-only helpers imported only by API routes. Optional. \\
\texttt{src/components/BuheraTerminal.js} & Two edits: import + register + optional renderer. \\
\bottomrule
\end{tabularx}

\subsection{Minimal skeleton (pure client, plain text output)}

The smallest working module. Create
\texttt{src/lib/modules/greeter-module.js}:

\begin{lstlisting}
export const greeterModule = {
  id: "greeter",

  describe() {
    return {
      id: "greeter",
      description: "greets whatever you send it",
      instructions: [
        'dispatch("greeter", "world")',
        'dispatch("greeter", { name: "Alice" })',
      ],
    };
  },

  async execute(instruction, _actBudget = 1) {
    const name = typeof instruction === "string"
      ? instruction
      : instruction?.name ?? "stranger";

    return {
      ok: true,
      output_delta: {
        kind: "text",
        lines: [`hello, ${name}`],
      },
      residue: 1,
      completed: true,
    };
  },

  outputCell(_instruction) {
    return { kind: "greeter_cell" };
  },
};
\end{lstlisting}

Then in \texttt{BuheraTerminal.js}:

\begin{lstlisting}
import { greeterModule } from "@/lib/modules/greeter-module";
// ...
register(greeterModule);
\end{lstlisting}

That's it. \texttt{next build} passes; \texttt{:modules} shows
\texttt{greeter}; \texttt{dispatch("greeter", "buhera")} returns
\texttt{hello, buhera}.

% =====================================================================
\section{Recipes by Module Type}
\label{sec:recipes}
% =====================================================================

Three concrete construction recipes, each corresponding to a common
class of scientific/domain module. Follow the recipe that matches the
module's shape.

\subsection{Recipe A: A pure computational module (no network)}

Use when the module runs an algorithm on inputs and returns a result.
Examples: a sequence-alignment module, a numerical-integration module,
a chemistry-formula parser.

\begin{recipe}[Pure computational]
\label{recipe:pure}
\begin{enumerate}[leftmargin=1.4em,itemsep=1pt]
\item Put the algorithm in \texttt{src/lib/<domain>/} as one or more
plain \texttt{.js} or \texttt{.ts} files. Test the algorithm in
isolation; the trait wrapper should be trivial.
\item Create \texttt{src/lib/modules/<id>-module.js} exporting a trait
object that imports the algorithm and wraps it.
\item Accept both a bare string (as \texttt{"demo"} or a domain-natural
form) and a structured object as instruction.
\item Return an ActResult with a domain-specific \texttt{output\_delta.kind}
and enough structured data for a future renderer.
\item Add a renderer if the output benefits from visual formatting.
\end{enumerate}
\end{recipe}

Lavoisier is the reference for this recipe. Look at
\texttt{src/lib/modules/lavoisier-module.js} for the trait wrapper and
\texttt{src/lib/lavoisier/experiment/virtualinstrument.js} for the
underlying algorithm.

\subsection{Recipe B: A module that calls an LLM}

Use when the module needs a language model in the loop. Examples: a
paper-summariser, a code-explainer, a natural-language SQL translator.

\begin{recipe}[LLM-backed]
\label{recipe:llm}
\begin{enumerate}[leftmargin=1.4em,itemsep=1pt]
\item Create a server route \texttt{src/pages/api/<id>.js} that accepts
POST, validates its inputs, calls \texttt{chatCascade} from
\texttt{src/lib/server/llm-cascade}, and returns JSON with
\texttt{ok/content/provider/model}.
\item Design the system prompt carefully. If the output must be
structured, define a JSON schema and pass it to \texttt{chatCascade} as
\texttt{jsonSchema}. Gemini and OpenAI honour it natively; HuggingFace
is skipped automatically when a schema is required.
\item Create \texttt{src/lib/modules/<id>-module.js} whose \texttt{execute}
POSTs to \texttt{/api/<id>} and translates the response into an
ActResult. Surface the provider used in the \texttt{output\_delta} so
the renderer can display it.
\item Renderer shows the LLM's response as either free text or a
structured card. Include the provider name as a small header.
\item In the module's \texttt{describe}, note that an LLM key is
required.
\end{enumerate}
\end{recipe}

Zangalewa is the reference for this recipe. See
\texttt{src/pages/api/extract.js} for the server route and
\texttt{src/lib/modules/zangalewa-module.js} for the client adapter.

\subsection{Recipe C: A module that reads history}

Use when the module's job is to observe the running system and produce
derived information. Examples: a session summariser, a token-usage
counter, a contradiction detector.

\begin{recipe}[Meta / history-reading]
\label{recipe:meta}
\begin{enumerate}[leftmargin=1.4em,itemsep=1pt]
\item Import \texttt{getAuditLog} from \texttt{@/lib/modules/registry}.
Do not import \texttt{onDispatch}; the hook mechanism is for the
terminal's wiring layer, not for modules.
\item In \texttt{execute}, read the audit log, filter or aggregate as
needed, and return the derived information.
\item If the module maintains its own state across dispatches (as
purpose-carry does with its Session), keep the state in a
module-private singleton and expose \texttt{reset}-style instructions
so callers can wipe it.
\item Consider whether the module's dispatches should themselves feed
back into any state the module reads. Purpose-carry avoids this by
skipping self-dispatches in the audit-log feeder. A new meta module
must make an equivalent decision explicitly.
\end{enumerate}
\end{recipe}

Purpose-carry is the reference for this recipe. See
\texttt{src/lib/modules/purpose-carry-module.js} for the module and
\texttt{src/components/BuheraTerminal.js} for how the audit-log hook
is wired (only the terminal registers the hook; the module does not).

% =====================================================================
\section{Instruction Design}
\label{sec:instruction-design}
% =====================================================================

The most important design decision for a new module is the shape of its
instruction. A well-designed instruction is small, orthogonal, and
covers three cases: a demo, a common single-value form, and a full
structured form.

\subsection{The three canonical shapes}

Every existing module accepts at least these:

\begin{enumerate}[nosep,leftmargin=1.4em]
\item \texttt{"demo"} or \texttt{""} --- runs a canned demonstration.
Zero-argument invocation is discoverable and testable.
\item A bare string --- the module interprets it in a domain-natural
way (a query, a source snippet, a name, a URL). This is the form
turbulance scripts most commonly emit.
\item A structured object \texttt{\{ kind: "action", ...args \}} ---
the fully explicit form used by callers who need control over every
parameter.
\end{enumerate}

Modules SHOULD accept all three. Modules MUST accept at least one.

\subsection{The \texttt{kind} discriminator}

When a module has multiple actions, the structured form MUST include a
\texttt{kind} field:

\begin{lstlisting}
dispatch("my-module", { kind: "search", query: "..." })
dispatch("my-module", { kind: "index",  document: "..." })
dispatch("my-module", { kind: "stats" })
dispatch("my-module", { kind: "reset" })
\end{lstlisting}

Purpose-carry uses this pattern (\texttt{carry}, \texttt{add},
\texttt{stats}, \texttt{reset}). It is the recommended pattern for any
module with more than one action.

\subsection{Backwards-compatible evolution}

Once shipped, an instruction shape is a contract. Evolving it safely:

\begin{itemize}[nosep,leftmargin=1.4em]
\item Adding a new optional field is safe.
\item Adding a new \texttt{kind} is safe.
\item Renaming a required field is a breaking change.
\item Changing the meaning of \texttt{"demo"} is a breaking change if
users script against it.
\end{itemize}

A module's tests should exercise every documented instruction shape so
that regressions in the shape are caught before shipping.

% =====================================================================
\section{Integrating with the Existing Federation}
\label{sec:integration}
% =====================================================================

New modules should compose with, not replace, existing ones. Three
common composition patterns.

\subsection{Reading from vaHera}

If a new module benefits from a semantic memory (retrieval by S-distance
against stored notes), it can read from vaHera's live kernel. Import
\texttt{getKernel} from \texttt{src/lib/modules/vahera-module.js} (which
exports it) and use \texttt{kernel.proximity(target, k)} or
\texttt{kernel.surgicalRetrieve(target, k)}. The graffiti kernel-search
catalyst is the reference implementation.

\subsection{Feeding purpose-carry}

Every dispatched act is automatically τ-extracted and fed into
purpose-carry's session by the terminal's post-dispatch hook. A new
module does not do anything special to appear in the carry --- it
happens automatically as long as the instruction contains extractable
terms.

If the module's instruction is inherently structured (say, a large
config object with no natural text content), consider including a
short descriptive string field in the instruction so the τ extractor
has something to work with. Example: lavoisier's config accepts a
\texttt{description} field alongside the numeric parameters.

\subsection{Being called by another module}

The registry's \texttt{dispatch} is callable from inside any module's
\texttt{execute}. This lets a module compose others. Caveats:

\begin{itemize}[nosep,leftmargin=1.4em]
\item Cross-module dispatch appears in the audit log with the inner
call's module id, not the outer's.
\item Deep chains are hard to debug. If a module calls three others in
sequence, the audit log will show four acts and a caller has to trace
the causal chain.
\item Circular dispatch (module A calls module B calls module A) will
loop forever. The registry does not detect this.
\end{itemize}

Use cross-module dispatch when the composition is genuinely one-way and
narrow: an interceptor calling a solver, a meta module calling a
solver, a synthesis module calling a search module.

% =====================================================================
\section{Testing a New Module}
\label{sec:testing}
% =====================================================================

Four tests every new module should pass before it lands in
\texttt{main}.

\subsection{Build}

\begin{lstlisting}[style=shell]
cd long-grass
npm run build
\end{lstlisting}

Must complete without errors. Warnings about caniuse-lite are cosmetic
and can be ignored.

\subsection{Lint}

\begin{lstlisting}[style=shell]
cd long-grass
npx next lint --dir src --max-warnings 999
\end{lstlisting}

Must report zero errors.

\subsection{Dispatch smoke test}

Boot \texttt{npm run dev}, open the terminal, and dispatch to the new
module. It MUST:

\begin{itemize}[nosep,leftmargin=1.4em]
\item Appear in \texttt{:modules} with the expected description.
\item Return an ActResult with \texttt{ok: true} for its \texttt{demo}
form.
\item Return an ActResult with \texttt{ok: false} and a specific error
for known bad input (empty string when a string is required, missing
required field, etc.).
\item Log to \texttt{:audit} with the correct \texttt{module\_id}.
\end{itemize}

\subsection{Federation smoke test}

Confirm that every other module still works after adding the new one.
\texttt{:modules} shows the correct count. A dispatch to
\texttt{purpose-carry} still reads the audit log correctly. If the new
module returns a domain-specific \texttt{output\_delta.kind}, the
renderer correctly renders it.

% =====================================================================
\part{Reference}
% =====================================================================

% =====================================================================
\section{Complete Example: A Sequence-Alignment Module}
\label{sec:example}
% =====================================================================

A fully worked example following Recipe A. This module computes a
Needleman-Wunsch global alignment between two DNA sequences and returns
the aligned strings, the score, and the CIGAR string.

\subsection{The algorithm}

Put in \texttt{src/lib/bioinformatics/needleman-wunsch.js}:

\begin{lstlisting}
export function align(seq1, seq2, opts = {}) {
  const match = opts.match ?? 1;
  const mismatch = opts.mismatch ?? -1;
  const gap = opts.gap ?? -1;

  const n = seq1.length, m = seq2.length;
  const dp = Array.from({length: n+1}, () => new Int32Array(m+1));

  for (let i = 0; i <= n; i++) dp[i][0] = i * gap;
  for (let j = 0; j <= m; j++) dp[0][j] = j * gap;

  for (let i = 1; i <= n; i++) {
    for (let j = 1; j <= m; j++) {
      const s = seq1[i-1] === seq2[j-1] ? match : mismatch;
      dp[i][j] = Math.max(
        dp[i-1][j-1] + s,
        dp[i-1][j] + gap,
        dp[i][j-1] + gap
      );
    }
  }

  // Traceback
  let a1 = "", a2 = "", cigar = "";
  let i = n, j = m;
  while (i > 0 || j > 0) {
    if (i > 0 && j > 0 &&
        dp[i][j] === dp[i-1][j-1] +
          (seq1[i-1] === seq2[j-1] ? match : mismatch)) {
      a1 = seq1[i-1] + a1;
      a2 = seq2[j-1] + a2;
      cigar = "M" + cigar;
      i--; j--;
    } else if (i > 0 && dp[i][j] === dp[i-1][j] + gap) {
      a1 = seq1[i-1] + a1;
      a2 = "-" + a2;
      cigar = "D" + cigar;
      i--;
    } else {
      a1 = "-" + a1;
      a2 = seq2[j-1] + a2;
      cigar = "I" + cigar;
      j--;
    }
  }

  return { seq1: a1, seq2: a2, score: dp[n][m], cigar };
}
\end{lstlisting}

\subsection{The module}

Put in \texttt{src/lib/modules/align-module.js}:

\begin{lstlisting}
import { align } from "@/lib/bioinformatics/needleman-wunsch";

const DEMO_INPUT = { seq1: "GATTACA", seq2: "GCATGCU" };

function resolveInput(instruction) {
  if (instruction == null || instruction === "" || instruction === "demo") {
    return DEMO_INPUT;
  }
  if (typeof instruction === "object" && instruction.kind === "align") {
    return {
      seq1: String(instruction.seq1 ?? "").toUpperCase(),
      seq2: String(instruction.seq2 ?? "").toUpperCase(),
      opts: instruction.opts || {},
    };
  }
  return null;
}

export const alignModule = {
  id: "align",

  describe() {
    return {
      id: "align",
      description:
        "global pairwise sequence alignment (Needleman-Wunsch). " +
        "Returns aligned sequences, score, and CIGAR.",
      instructions: [
        'dispatch("align", "demo")',
        'dispatch("align", { kind: "align", seq1: "ACGT", seq2: "AGGT" })',
      ],
    };
  },

  async execute(instruction, _actBudget = 1) {
    const input = resolveInput(instruction);
    if (!input) {
      return {
        ok: false,
        output_delta: {
          kind: "text",
          lines: [
            'align: instruction must be "demo" or ' +
              '{ kind: "align", seq1, seq2 }',
          ],
        },
        residue: 0,
        completed: true,
        error: "invalid instruction",
      };
    }
    if (!input.seq1 || !input.seq2) {
      return {
        ok: false,
        output_delta: {
          kind: "text",
          lines: ["align: seq1 and seq2 must both be non-empty"],
        },
        residue: 0,
        completed: true,
        error: "empty sequences",
      };
    }

    try {
      const result = align(input.seq1, input.seq2, input.opts);
      return {
        ok: true,
        output_delta: {
          kind: "align_result",
          seq1_aligned: result.seq1,
          seq2_aligned: result.seq2,
          score: result.score,
          cigar: result.cigar,
        },
        residue: Math.max(0, result.score),
        completed: true,
      };
    } catch (err) {
      return {
        ok: false,
        output_delta: {
          kind: "text",
          lines: [`align error: ${err.message || String(err)}`],
        },
        residue: 0,
        completed: true,
        error: err.message || String(err),
      };
    }
  },

  outputCell(_instruction) {
    return { kind: "align_cell" };
  },
};
\end{lstlisting}

\subsection{The renderer}

In \texttt{BuheraTerminal.js}, add a component near the other artifact
renderers:

\begin{lstlisting}
function ArtifactAlign({ seq1_aligned, seq2_aligned, score, cigar }) {
  return (
    <div className="text-gray-300 font-mono text-sm">
      <div className="text-xs text-gray-500 mb-1">
        score: <span className="text-white">{score}</span>
        {"  ·  "}CIGAR: <span className="text-white">{cigar}</span>
      </div>
      <pre className="whitespace-pre">{seq1_aligned}</pre>
      <pre className="whitespace-pre">{seq2_aligned}</pre>
    </div>
  );
}
\end{lstlisting}

And add a case in the artifact switch:

\begin{lstlisting}
case "align_result":
  return <ArtifactAlign
    seq1_aligned={result.seq1_aligned}
    seq2_aligned={result.seq2_aligned}
    score={result.score}
    cigar={result.cigar}
  />;
\end{lstlisting}

\subsection{Registration}

At the top of \texttt{BuheraTerminal.js}:

\begin{lstlisting}
import { alignModule } from "@/lib/modules/align-module";
\end{lstlisting}

In the mount effect:

\begin{lstlisting}
register(alignModule);
\end{lstlisting}

\subsection{Verification}

\begin{lstlisting}[style=shell]
npm run build     # should pass
npm run dev
# in the browser terminal:
:modules          # should list "align"
dispatch("align", "demo")
# expected: aligned GATTACA and GCATGCU, some negative score, CIGAR
dispatch("align", { kind: "align", seq1: "ACGT", seq2: "AGGT" })
# expected: alignment showing one mismatch
:audit            # two entries with module_id: "align"
\end{lstlisting}

That is a complete, working, testable module. Every scientific module
follows this same shape at a higher level of domain complexity.

% =====================================================================
\section{Common Pitfalls}
\label{sec:pitfalls}
% =====================================================================

Nine mistakes I have seen (or made) that break a module in ways whose
error messages are unhelpful.

\begin{enumerate}[leftmargin=1.6em,itemsep=2pt]
\item \textbf{Forgetting to make \texttt{execute} async.} The registry
awaits the return value; a non-async function that returns a plain
object also happens to work, but the moment the function does any
\texttt{await} inside it, forgetting \texttt{async} is a syntax error.
Always declare \texttt{async execute(...)}.

\item \textbf{Throwing instead of returning \texttt{ok:false}.} The
registry catches exceptions, but the resulting error is generic and the
audit-log entry loses domain context. Wrap the module's work in
\texttt{try/catch} and return proper ActResults.

\item \textbf{Naming an \texttt{output\_delta.kind} that clashes with an
existing renderer.} If a new module returns \texttt{kind: "text"} with
non-\texttt{lines} fields, the text renderer will silently do the wrong
thing. Namespace kinds with the module id.

\item \textbf{Importing server-side code into the client bundle.} A
module that imports \texttt{fs}, \texttt{child\_process}, or any
Node-only package into \texttt{src/lib/modules/} will break the build.
Server code lives under \texttt{src/pages/api/} or
\texttt{src/lib/server/}; nothing in \texttt{src/lib/modules/} may
transitively import Node built-ins.

\item \textbf{Storing state on the module trait object itself.} The
module object is a singleton, but so are its fields. Mutating
\texttt{myModule.someState = ...} inside \texttt{execute} creates
subtle bugs. Use a module-private variable (as vahera-module and
purpose-carry-module do with their \texttt{\_kernel} and \texttt{\_session}
singletons) and export a \texttt{reset} function if needed.

\item \textbf{Blocking the render thread with a long synchronous
computation.} The module runs in the browser's main thread. A CPU-bound
algorithm running for seconds will freeze the terminal. Break large
computations across microtasks (\texttt{await new Promise(r => setTimeout(r, 0))})
or offload to a Web Worker.

\item \textbf{Depending on module load order.} Modules are registered
in the mount effect in a fixed order, but a module MUST NOT rely on
another module being registered before it. Cross-module dispatch reads
the registry at call time, not at load time, so ordering only matters
for early-boot calls; even those should use defensive lookups.

\item \textbf{Assuming the audit log persists.} It resets on page
reload. A meta module that needs long-term memory must persist its own
state (\texttt{localStorage}, \texttt{IndexedDB}, or a snapshot API).

\item \textbf{Skipping the outputCell stub.} The registry's
\texttt{register()} does not check for \texttt{outputCell}, but the
future scheduler gate will. Add the two-line stub now to avoid the
churn later.
\end{enumerate}

% =====================================================================
\section{Where to Look in the Existing Codebase}
\label{sec:references}
% =====================================================================

Every claim in this document corresponds to code in the shipped
Buhera. When in doubt, read the referenced file.

\begin{center}
\begin{tabularx}{\linewidth}{lX}
\toprule
\textbf{Topic} & \textbf{File} \\
\midrule
Module trait definition & \texttt{src/lib/modules/registry.js} \\
Reference: pure computational module & \texttt{src/lib/modules/lavoisier-module.js} \\
Reference: DSL-executing module & \texttt{src/lib/modules/vahera-module.js} \\
Reference: DSL search module & \texttt{src/lib/modules/graffiti-module.js} \\
Reference: LLM-backed module & \texttt{src/lib/modules/zangalewa-module.js} \\
Reference: meta / history-reading & \texttt{src/lib/modules/purpose-carry-module.js} \\
Reference: server-side federation & \texttt{src/lib/modules/purpose-module.js} \\
LLM cascade helper & \texttt{src/lib/server/llm-cascade.js} \\
Sample server route (LLM) & \texttt{src/pages/api/extract.js} \\
Sample server route (JS federation) & \texttt{src/pages/api/purpose-federation.js} \\
Provider introspection & \texttt{src/pages/api/providers.js} \\
Terminal registration + hooks & \texttt{src/components/BuheraTerminal.js} \\
τ (term extraction) & \texttt{src/lib/purpose-terms.js} \\
Cost estimation & \texttt{src/lib/purpose-cost.js} \\
\bottomrule
\end{tabularx}
\end{center}

% =====================================================================
\section{Checklist Before Shipping a New Module}
\label{sec:checklist}
% =====================================================================

Print this. Tick every line.

\begin{itemize}[leftmargin=1.6em,itemsep=1pt]
\item[$\square$] Module id chosen; lowercase, kebab-case if compound.
\item[$\square$] File \texttt{src/lib/modules/<id>-module.js} created and exports the trait object.
\item[$\square$] \texttt{describe()} returns \texttt{\{id, description, instructions\}}.
\item[$\square$] \texttt{execute()} is \texttt{async}, wrapped in \texttt{try/catch}, returns a valid ActResult in every branch.
\item[$\square$] Instruction handling covers \texttt{"demo"}, bare string, structured object.
\item[$\square$] Every ActResult has \texttt{ok}, \texttt{output\_delta}, \texttt{residue}, \texttt{completed}, and \texttt{error} iff \texttt{ok:false}.
\item[$\square$] \texttt{output\_delta.kind} is namespaced (starts with the module id).
\item[$\square$] \texttt{outputCell} stub present.
\item[$\square$] If LLM-backed: server route created; uses \texttt{chatCascade}; adapter posts to route; no direct provider calls in client bundle.
\item[$\square$] If it has a domain-specific output: renderer written; case added to the artifact switch.
\item[$\square$] Import and \texttt{register(myModule)} added to \texttt{BuheraTerminal.js} mount effect.
\item[$\square$] \texttt{npm run build} passes.
\item[$\square$] \texttt{npx next lint --dir src} passes.
\item[$\square$] \texttt{:modules} lists the new module with correct description.
\item[$\square$] \texttt{dispatch("<id>", "demo")} returns \texttt{ok:true} in the terminal.
\item[$\square$] \texttt{:audit} shows the dispatched act.
\item[$\square$] At least one other module (e.g. \texttt{purpose-carry}) still works after registration.
\item[$\square$] \texttt{tutorial.md} updated with at least one worked example of the new module.
\end{itemize}

% =====================================================================
\section{Closing}
\label{sec:closing}
% =====================================================================

The federation is designed to grow by addition. Every architectural
decision in Buhera exists to make it possible for a new module to land
without requiring anyone to understand the whole system. Read the trait
in \S\ref{sec:registry}, pick the recipe in \S\ref{sec:recipes} that
matches the module's shape, copy the skeleton in \S\ref{sec:skeleton},
consult the worked example in \S\ref{sec:example}, and run through the
checklist in \S\ref{sec:checklist}. What lands is a module that speaks
the same protocol every existing module speaks, appears in
\texttt{:modules}, dispatches from turbulance, contributes to the audit
log, and interoperates with purpose-carry, graffiti, zangalewa, and the
rest without any of them knowing about it.

Every scientific tool that the group produces --- protein structure,
molecular docking, chemical reaction networks, imaging pipelines,
statistical analyses, phylogenetics, whatever --- can become a Buhera
module by following this document. That is what the federation is for.

\end{document}
